
The modern software engineering landscape is currently undergoing a profound paradigm shift driven by the relentless acceleration of deployment lifecycles and the exponential growth of codebase complexities. This acceleration has been heavily catalyzed by the advent and widespread adoption of \textbf{Large Language Models (LLMs)} and \textbf{AI-assisted programming agents}, which have dramatically increased developer velocity. However, this fundamental shift in how software is authored introduces unprecedented, systemic risks to long-term software health. While AI-generated code enhances functional delivery speeds, it fundamentally alters the spatiotemporal evolution of software quality, often introducing subtle structural vulnerabilities, architectural inconsistencies, and latent technical debt that evade immediate detection\parencite{human_written_vs_ai_generated_code_2025}. Consequently, traditional software quality assurance mechanisms, which rely heavily on manual human review and static, point-in-time rule-checking, are rapidly becoming operational bottlenecks.

Software quality is a multifaceted construct, defined theoretically by international frameworks such as the \textbf{ISO/IEC 25010 standard}, which delineates non-functional characteristics including Maintainability, Security, Performance Efficiency, and Reliability \parencite{isoiec25010_quality_model_emergentmind}. Operationally, these abstract characteristics are measured utilizing static code analysis tools. \textbf{SonarQube} has established itself as the industry standard for continuous code quality inspection, utilizing an extensive taxonomy of deterministic rules to detect bugs, code smells, and security vulnerabilities across varied codebases \parencite{sonarqube_metrics_definition}. Yet, conventional static analysis remains fundamentally retrospective and static. It evaluates the codebase at a discrete point in time, identifying existing flaws but comprehensively failing to anticipate how current architectural decisions, coupled with the continuous influx of AI-generated snippets, will degrade maintainability or security in the future.

To address the limitations of retrospective analysis and to accommodate the profound nuances of AI-augmented software development, there is an imperative need for intelligent, predictive analytics platforms. The proposed research project, \textbf{Deltx}, aims to architect an expert-level static analysis and predictive analytics platform built atop the SonarQube engine. By integrating advanced machine learning techniques for stylometric AI-code detection and utilizing state-of-the-art Transformer-based time-series forecasting architectures, Deltx seeks to transform static, isolated metrics into proactive, explainable quality trajectories. This proposal details the comprehensive architectural design, mathematical modeling, and rigorous data strategy required to operationalize this predictive software quality framework.

\section{Problem Statement}

Despite the ubiquity of continuous integration and continuous deployment (\textbf{CI/CD}) pipelines, the software engineering domain faces a critical research gap in predicting and managing long-term software quality, a challenge specifically exacerbated by the opaque integration of AI-generated code \parencite{developers_self_declaration_ai_generated_code}. This project addresses four interrelated, foundational problems within this domain:

\subsubsection*{Issue 1: Invisibility of AI-Generated Structural Decay}

Current static analysis tools process all source code uniformly, failing entirely to distinguish between human-authored logic and AI-generated code. Artificial intelligence models, trained on vast but generalized, context-agnostic corpora, often generate code that is functionally correct in isolation but structurally suboptimal when integrated into broader enterprise architectures \parencite{human_written_vs_ai_generated_code_2025}. Recent empirical analyses indicate that AI-generated pull requests often contain up to 1.7 times more structural issues than human-authored equivalents, leading to substantially higher issue densities and rapid technical debt accumulation \parencite{coderabbit_ai_vs_human_code_generation_report}. Without the algorithmic capability to accurately identify AI-generated code, organizations cannot apply targeted verification, nor can predictive models accurately weigh the specific decay risks associated with synthetic authorship.

\subsubsection*{Issue 2: Semantic Gap in Software Quality Metrics}

Tools like SonarQube generate thousands of granular, highly specific rule-based metric outputs, tracking variables such as cyclomatic complexity, raw code smells, and arbitrary technical debt ratios \parencite{sonarqube_metrics_definition}. However, there is a distinct semantic gap between these raw, disjointed metrics and high-level business risk frameworks, specifically the ISO/IEC 25010 standard \parencite{isoiec25010_quality_model_emergentmind}. Translating raw, unweighted metrics into normalized, mathematically sound quality dimensions—such as Correctness, Maintainability, Security, and Efficiency—remains an unsolved architectural challenge. Existing aggregations often fail to handle conflicting signals, meaning a massive volume of trivial formatting warnings can mathematically overshadow a single, catastrophic security vulnerability \parencite{software_quality_metrics_aggregation_industry}.

\subsubsection*{Issue 3: Retrospective Nature of Static Analysis}

Platforms like SonarQube provide valuable point-in-time snapshots of codebase health but critically lack the capability to perform multivariate time-series forecasting to predict how code quality will evolve across future commits \parencite{fault_prediction_software_metrics_sonarqube_rules}. Without temporal forecasting, engineering teams are relegated to a reactive posture, addressing technical debt only after it breaches critical operational thresholds, rather than proactively managing structural health based on predicted degradation curves \parencite{fault_prediction_software_metrics_sonarqube_rules}.

\subsubsection*{Issue 4: Lack of Interpretability in Predictive Metrics}

There is a profound lack of interpretability in predictive software metrics. The application of deep learning algorithms in software engineering often results in opaque, ``black-box'' predictions. If a predictive system forecasts a severe decline in software security over the next thirty days, it must concurrently provide human-readable, actionable insights explaining precisely why the forecast was made \parencite{perspective_on_xai_methods_shap_and_lime}. Without granular feature attribution—identifying the specific historical metrics, code churn events, or AI-code influxes driving the prediction—developers cannot trust or act upon the algorithmic insights, rendering the predictive platform practically useless \parencite{comparative_analysis_lime_shap_interpreters_diabetes}.

\vspace{0.5cm}

By systematically addressing these four challenges, Deltx will advance software engineering theory by combining stylometric AI detection, graph-theoretic bug weighting, advanced time-series forecasting, and Explainable AI into a unified, rigorous framework.